% ===================================================================
%  ТАБЛИЦА № 14 — Улица. — Торговцы.
%  Traduction 2026 d'apres texto/io, controlee sur texto/fr, texto/en
%  et texto/es. Consigne : voir texto/ru/10-tablica-01.tex.
% ===================================================================

%%K t14-apar-1 apar
\VUsaut{4.0mm}
\VUtitre{15.0pt}{60}{ТАБЛИЦА № 14}
\VUsaut{3.0mm}

%%K t14-tit-1 sub
\VUpk{11.4pt}{40}{Улица. --- Торговцы.}
\VUsaut{3.0mm}

%%K t14-01-1 p
1. --- Мой друг Иван, который живёт в деревне, приехал провести три дня у
моих. \VUgras{До} сих пор ему ни разу не случалось видеть большого города,
так что мы целыми днями гуляем, и я объясняю ему всё, чего он не знает.

%%K t14-02-1 p
2. --- Сперва мы пошли смотреть \VUgras{обсерваторию} \textsuperscript{(1)},
которая очень его заняла. Возвращаясь по \VUgras{улице} \textsuperscript{(3)}, где
находятся \VUgras{турецкие бани} \textsuperscript{(2)}, мы встретили
\VUgras{похороны} \textsuperscript{(4)}. Во главе \VUgras{погребального шествия}
шло \VUgras{духовенство}, впереди \VUgras{катафалка}, на который поставили
\VUgras{гроб}. Много друзей шло за семьёй в \VUgras{трауре}.

%%K t14-02-2 p
Когда шествие прошло, мы перешли улицу перед большой \VUgras{пивной}
\textsuperscript{(5)}, где много \VUgras{посетителей} пило \VUgras{пиво} на
террасе под застеклённым \VUgras{навесом} \textsuperscript{(6)}.

%%K t14-03-1 p
3. --- Ивана озадачил \VUgras{герб}, помещённый между лавкой
\VUgras{торговца винами и водками} \textsuperscript{(7)} и лавкой
\VUgras{продавца нот} \textsuperscript{(10)}. Он спросил меня, что это значит; я
сказал ему, что это английское \VUgras{консульство} \textsuperscript{(8)}, и
показал \VUgras{консула} \textsuperscript{(9)}, стоявшего на балконе.

%%K t14-tit-2 sub
\VUpk{10.2pt}{40}{Разглядывая витрины.}
\VUsaut{3.0mm}

%%K t14-04-1 p
4. --- Разумеется, мы остановились перед выставкой
\VUgras{музыкальных инструментов}, \VUgras{фисгармоний}, \VUgras{органов} и
\VUgras{роялей}; мы разглядывали иллюстрированные обложки \VUgras{нот} и
фортепианных \VUgras{пьес} и любовались золочёной деревянной \VUgras{лирой},
служившей вывеской.

%%K t14-05-1 p
5. --- Облака пыли, поднятые \VUgras{дворниками} \textsuperscript{(11)} и
\VUgras{подметальной машиной} \textsuperscript{(12)}, заставили нас поскорее
пройти мимо прекрасных \VUgras{мебельных гарнитуров} у \VUgras{мебельщика}
\textsuperscript{(13)} и мимо выставки \VUgras{печей} и \VUgras{ламп} в
\VUgras{скобяной лавке} \textsuperscript{(14)}.

%%K t14-06-1 p
6. --- В этом самом месте нам и вправду пришлось сойти с тротуара, потому
что он был загромождён тюками, которые снимали с \VUgras{мебельного фургона}
\textsuperscript{(16)}, чтобы внести в только что снятую \VUgras{лавку}
\textsuperscript{(15)}.

%%K t14-06-2 p
Это помешало нам увидеть прекрасные экипажи у \VUgras{каретника}
\textsuperscript{(17)}, но не помешало остановиться и посмотреть, как
\VUgras{упаковщик} \textsuperscript{(19)} сколачивает ящик для своего соседа,
\VUgras{торговца велосипедами и автомобилями} \textsuperscript{(18)}. Потом, так
как было уже довольно поздно, мы вернулись домой.

%%K t14-tit-3 sub
\VUpk{10.2pt}{40}{Прогулка по городу.}
\VUsaut{3.0mm}

%%K t14-07-1 p
7. --- На другой день мать предложила, чтобы Иван снялся, и попросила меня
проводить его к \VUgras{фотографу} \textsuperscript{(20)}. После завтрака мы
пошли в \VUgras{ателье} художника, где пришлось порядочно подождать. Чтобы
скоротать время, мы рассматривали \VUgras{портреты} \textsuperscript{(22)},
висевшие на стене.

%%K t14-07-2 p
Когда пришла наша очередь, дело вышло недолгое, и, едва мой друг кончил
позировать перед \VUgras{аппаратом} \textsuperscript{(21)}, мы спустились.

%%K t14-08-1 p
8. --- На первом этаже мы увидели через открытую дверь
\VUgras{парикмахерской} \textsuperscript{(23)}, как подмастерье стриг клиента.

%%K t14-08-2 p
Снова выйдя на улицу, мы прошли мимо \VUgras{табачной лавки} \textsuperscript{(24)}.
Ивана так позабавили красный \VUgras{фонарь} \textsuperscript{(25)} и
\VUgras{большая трубка}, служившая \VUgras{вывеской} \textsuperscript{(26)}, что он
налетел на двух стариков, беседовавших за \VUgras{понюшкой табаку}
\textsuperscript{(27)}.

%%K t14-tit-4 sub
\VUpk{10.2pt}{40}{Кондитерская.}
\VUsaut{3.0mm}

%%K t14-09-1 p
9. --- \VUgras{Ассигнации} и \VUgras{монеты} всех стран, выставленные в
\VUgras{витрине} \VUgras{менялы} \textsuperscript{(29)}, задержали нас на минуту,
но так как был час полдника, мы, не мешкая долее, вошли в
\VUgras{кондитерскую} \textsuperscript{(31)}.

%%K t14-10-1 p
10. --- При виде всех \VUgras{лакомств}, какие были в лавке ---
\VUgras{пирожных, пряников, печенья} и \VUgras{конфет} всякого рода --- мы
насилу могли выбрать. В конце концов Иван остановился на
\VUgras{кремовых пирожных}, а я взял только маленький
\VUgras{вишнёвый тарт}, потому что от сладкого \VUgras{у меня болят зубы}, а
мне вовсе не хочется слишком часто ходить к \VUgras{зубному врачу}
\textsuperscript{(30)}.

%%K t14-tit-5 sub
\VUpk{10.2pt}{40}{Писание на машинке.}
\VUsaut{3.0mm}

%%K t14-11-1 p
11. --- Чтобы показать другу \VUgras{пишущую машинку}, о которой он слышал,
но которой не видел, мы зашли в \VUgras{контору} \textsuperscript{(32)} моего
\VUgras{двоюродного брата} \textsuperscript{(34)}. Мы застали его диктующим письма
своей \VUgras{секретарше} \textsuperscript{(35)}, которая \VUgras{стенографистка}.
Едва письмо было кончено, \VUgras{машинист} \textsuperscript{(33)} переписывал его
на своей машинке. Не желая мешать работе родственника, мы пробыли у него
всего несколько минут.

%%K t14-12-1 p
12. --- Под конторой двоюродного брата живёт большой
\VUgras{часовщик и ювелир} \textsuperscript{(37)}, который вдобавок серебряных дел
мастер. В его великолепной лавке много \VUgras{стенных часов, каминных
часов, карманных часов}, \VUgras{цепочек} и дорогих \VUgras{драгоценностей}:
\VUgras{жемчужных ожерелий}, золотых \VUgras{браслетов}, \VUgras{серёг} и
\VUgras{колец} с \VUgras{драгоценными камнями} и \VUgras{бриллиантами}. Одна
из витрин отдана целиком \VUgras{серебру} и вмещает большое собрание
серебряных \VUgras{приборов}.

%%K t14-13-1 p
13. --- Заворачивая за угол улицы, я заметил, что \VUgras{дощечку}
\textsuperscript{(36)} рядом с \VUgras{номером} дома переменили. Я собирался
сказать это вслух, когда раздались весёлые звуки военного оркестра. Мы
пустились бежать к полку, не подумав, что он пройдёт мимо лавок
\VUgras{шляпника} \textsuperscript{(39)} и \VUgras{перчаточника} \textsuperscript{(40)} и
что мы были бы так же хорошо поставлены, чтобы смотреть на его прохождение,
если бы остались перед витриной \VUgras{оптика} \textsuperscript{(38)}, полной
\VUgras{пенсне, очков} и \VUgras{биноклей}.

%%K t14-tit-6 sub
\VUpk{10.2pt}{40}{Прохождение полка.}
\VUsaut{3.0mm}

%%K t14-14-1 p
14. --- Во главе \VUgras{полка} \textsuperscript{(41)} шёл \VUgras{тамбурмажор}
\textsuperscript{(42)}, за ним \VUgras{барабанщики} \textsuperscript{(43)},
\VUgras{горнисты} \textsuperscript{(44)} и весь \VUgras{оркестр}. \VUgras{Генерал}
\textsuperscript{(45)}, который был на площади со своим адъютантом, остановился
возле \VUgras{струи} \textsuperscript{(49)} \VUgras{фонтана} \textsuperscript{(48)},
чтобы смотреть \VUgras{парад}.

%%K t14-14-2 p
\VUgras{Штабные офицеры} \textsuperscript{(46)}, \VUgras{полковник} и
\VUgras{подполковник} салютовали ему шпагами. \VUgras{Майоры} шли во главе
своих \VUgras{батальонов} \textsuperscript{(47)}, и каждый \VUgras{капитан}
командовал своей \VUgras{ротой}, а \VUgras{поручики}, \VUgras{подпоручики} и
\VUgras{унтер-офицеры} занимали свои обычные места рядом со своими людьми.

%%K t14-15-1 p
15. --- Много прохожих собралось на \VUgras{перекрёстке} \textsuperscript{(50)};
лавочники --- \VUgras{зонтичник} \textsuperscript{(51)}, \VUgras{красильщик}
\textsuperscript{(53)}, \VUgras{оружейник} \textsuperscript{(54)} --- вышли посмотреть
на \VUgras{солдат}. Все экипажи --- \VUgras{извозчичьи пролётки}
\textsuperscript{(52)}, \VUgras{собственные экипажи} \textsuperscript{(57)}, а также
\VUgras{поливальная бочка} \textsuperscript{(56)} --- прижались к тротуару.

%%K t14-tit-7 sub
\VUpk{10.2pt}{40}{Сцены на улице.}
\VUsaut{3.0mm}

%%K t14-15-2 p
\VUgras{Коммивояжёр} \textsuperscript{(55)}, несший в руках свои
\VUgras{чемоданы с образцами}, живо перешёл улицу, а сидевшие на
\VUgras{империале} \textsuperscript{(59)} \VUgras{омнибуса} \textsuperscript{(58)}
оборачивались, чтобы полюбоваться зрелищем. Бедный \VUgras{уличный певец}
\textsuperscript{(60)} со своей \VUgras{шарманкой} \textsuperscript{(61)} должен был
прервать свою \VUgras{песню}, потому что шум барабанов заглушал его голос.

%%K t14-16-1 p
16. --- Когда полк поравнялся с \VUgras{суконной лавкой} \textsuperscript{(64)},
\VUgras{швеи} \textsuperscript{(63)} и \VUgras{портные} \textsuperscript{(62)} бросили
свои \VUgras{швейные машины} и работу, чтобы выглянуть в окно.
\VUgras{Лавочник} \textsuperscript{(65)}, который только что выставил новые
\VUgras{костюмы} \textsuperscript{(66)} в свою \VUgras{витрину}, должен был убрать
штуки \VUgras{сукна} \textsuperscript{(67)} и \VUgras{полотна}, разложенные на
тротуаре возле \VUgras{водостока} \textsuperscript{(68)}, боясь, что толпа их
опрокинет; даже старому \VUgras{коробейнику} \textsuperscript{(69)}, у которого
привратница торговала чулки, пришлось зайти в подворотню, чтобы докончить
продажу.

%%K t14-16-2 p
Мы проводили полк до казарм \VUgras{и}, очень довольные своим днём,
вернулись домой к обеду.

%%K t14-tit-8 sub
\VUpk{10.2pt}{40}{Разные ремёсла и занятия.}
\VUsaut{3.0mm}

%%K t14-17-1 p
17. --- На другой день отец предложил свести нас в типографию местной
газеты. Мы приняли с восторгом.

%%K t14-17-2 p
\VUgras{Типограф} --- он же и \VUgras{книгопродавец} \textsuperscript{(70)},
\VUgras{издатель}, \VUgras{писчебумажник} и \VUgras{переплётчик}. Он поместил
на первом этаже \VUgras{редакцию} \textsuperscript{(71)} газеты, а также
\VUgras{гравёрную}, где работают \VUgras{литографы} \textsuperscript{(72)} и
\VUgras{гравёры по меди} \textsuperscript{(73)}, и, наконец, наборную, где сидят
\VUgras{наборщики} \textsuperscript{(74)}. Собственно \VUgras{машинный зал}
\textsuperscript{(75)}, \VUgras{печатные станки} и разные машины --- в нижнем
этаже.

%%K t14-tit-9 sub
\VUpk{10.2pt}{40}{Иван задаёт множество вопросов.}
\VUsaut{3.0mm}

%%K t14-18-1 p
18. --- Выйдя из типографии, Иван остановился в большом недоумении перед
стеклянным ящичком на столбе против \VUgras{галантерейной лавки}
\textsuperscript{(77)} и спросил, для чего он. Отец объяснил, что это
\VUgras{пожарный извещатель} \textsuperscript{(76)}. И правда, всё в городе было
для моего друга чудом: и простой \VUgras{питьевой фонтанчик} \textsuperscript{(78)},
и лёгкий \VUgras{развозной трицикл} \textsuperscript{(80)}, на котором ехал
\VUgras{рассыльный} \textsuperscript{(79)} в ливрее, и \VUgras{почтовый фургон}
\textsuperscript{(81)}.

%%K t14-19-1 p
19. --- Пока отец оставил нас на минуту, чтобы поговорить со
\VUgras{сборщиком податей} \textsuperscript{(83)}, который остановился побеседовать
с \VUgras{адвокатом} \textsuperscript{(82)} и \VUgras{нотариусом} \textsuperscript{(84)},
мы забавлялись, следя за движением на улице. Мимо нас проходил самый
разный народ: \VUgras{кондитерский мальчик} \textsuperscript{(85)}, нёсший в
корзине великолепный \VUgras{пирог}, шёл за \VUgras{старьёвщиком}
\textsuperscript{(86)}; \VUgras{фонарщик} \textsuperscript{(87)} беседовал с
\VUgras{газовщиком} \textsuperscript{(88)}; \VUgras{монахиня} \textsuperscript{(89)},
державшая за руку маленькую сиротку, возвращалась в свой монастырь.

%%K t14-tit-10 sub
\VUpk{10.2pt}{40}{На обратном пути.}
\VUsaut{3.0mm}

%%K t14-20-1 p
20. --- Как раз когда отец к нам вернулся, Иван расхохотался при виде
чумазых лиц и странного наряда двух \VUgras{трубочистов} \textsuperscript{(91)},
чёрных от сажи.

%%K t14-21-1 p
21. --- Я хотел купить матери растение у \VUgras{цветочницы} \textsuperscript{(92)},
но отец заметил, что его тяжело будет нести и что лучше взять
\VUgras{апельсинов}. Мы подошли к \VUgras{торговке с лотка} \textsuperscript{(94)},
и, когда \VUgras{гувернантка} \textsuperscript{(93)}, выбиравшая их своему
питомцу, кончила покупку, нас обслужили.

%%K t14-22-1 p
22. --- На обратном пути мы встретили нескольких знакомых: старую приятельницу
нашей семьи, опиравшуюся на руку своей \VUgras{компаньонки} \textsuperscript{(95)},
\VUgras{инженера} \textsuperscript{(96)} путей сообщения и одну из моих
двоюродных сестёр с её \VUgras{няней} \textsuperscript{(98)}.

%%K t14-22-2 p
Потом мы прошли мимо \VUgras{модистки} \textsuperscript{(97)} моей матери и двух
\VUgras{монахов} \textsuperscript{(99)}, приезжих в городе, которые подошли
спросить у отца дорогу на вокзал.
\VUsaut{6.0mm}
\VUfilet{22mm}
